% ============================================================
\section{Conclusions}
% ============================================================

The fundamental tension in spine biomechanics research --- between the precision of
laboratory optical capture and the ecological validity of ambulatory field measurement
--- has not been resolved by conventional \acs{IMU} systems.
Orientation integration accumulates drift, a single sensor on the thorax loses
segmental context, and multi-sensor arrays introduce complexity that precludes
long-term free-living deployment.

Differential strain-gauge shape reconstruction addresses these constraints at the
hardware level.
By encoding curvature directly rather than integrating angular velocity, the sensor
provides drift-free segmental resolution of the full thoracolumbar spine.
Anchoring the reconstruction at the sacrum with a magnetometer-free \acs{IMU} preserves
this accuracy in magnetically disturbed environments.

The evidence reviewed here spans three independent research groups across three
published peer-reviewed studies and one pre-publication in-vitro validation:

\begin{itemize}
	\item Geometric accuracy matching clinical-grade instruments at \ang{0.8}--\ang{1.1}
	      \acs{MAE} \cite{walkling2025}.
	\item Activity recognition at F1\,=\,0.90 (1-second window) across 11 daily
	      activities, with higher trunk-posture accuracy than a video baseline
	      \cite{walkling2025}.
	\item Quantification of real-world ergonomic improvements on an active construction
	      site where no alternative sensor would function \cite{sawicki2026}.
	\item Post-surgical \acs{ROM} measurement in a clinical pilot \cite{marx2023}.
\end{itemize}

Taken together, these results support positioning the FlexTail system as a capable
research instrument for any study in which (a) continuous ambulatory spine kinematics
are needed, (b) the measurement environment is magnetically disturbed or optically
obstructed, or (c) participant burden must be low enough to sustain multi-day or
multi-week wear.

Researchers considering the system for clinical trials, occupational cohort studies,
or sport biomechanics programmes are encouraged to request the raw-data specification
and the \emph{flexlib} documentation from MinkTec to assess suitability for their
specific measurement demands.
